\chapter{Najczęstsze błędy studentów}

\begin{summarybox}
Ten dodatek zbiera błędy, które najczęściej psują obliczenia, wykresy i raporty z
praktycznej mechaniki kwantowej. Większość z nich nie wynika z braku inteligencji,
lecz z braku kontroli jednostek, normalizacji, warunków brzegowych, składni
Mathematiki i interpretacji fizycznej wyniku.
\end{summarybox}

\section{Wykres w dżulach zamiast w eV}

Bardzo częsty błąd polega na policzeniu energii w dżulach i narysowaniu wykresu
bez konwersji na eV. Liczby są wtedy ekstremalnie małe, wykres wygląda dziwnie, a
student zaczyna poprawiać fizykę zamiast jednostek.

Jeżeli energia została policzona w SI, należy wykonać
\begin{equation}
    E_{\mathrm{eV}}=\frac{E_{\mathrm{J}}}{1.602176634\times10^{-19}}.
\end{equation}
Na osiach wykresu trzeba zawsze podać jednostkę. Opis osi „Energy” bez jednostki
jest niepełny.

\section{Szerokość w metrach zamiast w nm}

Jeśli szerokość studni wynosi \(L=10\,\mathrm{nm}\), to w jednostkach SI jest to
\begin{equation}
    L=10^{-8}\,\mathrm{m}.
\end{equation}
Podstawienie \(L=10\) do wzoru SI jest dramatycznym błędem. Ponieważ energia w
studni zależy od \(1/L^2\), błąd długości o czynnik \(10^9\) daje błąd energii o
czynnik \(10^{18}\).

W kodzie warto pisać:
\begin{verbatim}
nm = 10^-9;
L = 10 nm;
\end{verbatim}
zamiast wpisywać gołe liczby bez jednostkowego komentarza.

\section{Brak normalizacji funkcji falowej}

Funkcja falowa musi być znormalizowana:
\begin{equation}
    \int |\psi(x)|^2\,dx=1.
\end{equation}
Jeśli funkcja nie jest znormalizowana, wartości oczekiwane i prawdopodobieństwa
będą błędne. W Mathematice warto zawsze sprawdzić:
\begin{verbatim}
Integrate[Conjugate[psi[x]] psi[x], {x, xmin, xmax}]
\end{verbatim}
albo numerycznie:
\begin{verbatim}
NIntegrate[Abs[psi[x]]^2, {x, xmin, xmax}]
\end{verbatim}

\section{Mylenie \texorpdfstring{\(\psi\)}{psi} i \texorpdfstring{\(|\psi|^2\)}{|psi|^2}}

Funkcja \(\psi\) jest amplitudą, a \(|\psi|^2\) jest gęstością
prawdopodobieństwa. To nie są te same obiekty. Amplituda może być dodatnia,
ujemna albo zespolona. Gęstość prawdopodobieństwa jest nieujemna.

Na wykresach trzeba jasno opisywać, co jest rysowane. Jeśli rysujesz część
rzeczywistą funkcji falowej, podpisz \(\operatorname{Re}\psi\). Jeśli rysujesz
gęstość, podpisz \(|\psi|^2\). Nie należy opisywać obu jako „wave function”.

\section{Brak warunków brzegowych}

Równanie Schrödingera bez warunków brzegowych zwykle nie definiuje pełnego
problemu własnego. Dla studni nieskończonej trzeba narzucić znikanie funkcji
falowej na ścianach. Dla stanów związanych na całej osi trzeba wymagać
normalizowalności. Dla problemów rozproszeniowych trzeba określić falę padającą,
odbitą i transmitowaną.

Brak warunków brzegowych prowadzi do rozwiązań matematycznych, które nie są
fizycznym rozwiązaniem konkretnego zadania.

\section{Złe założenia w \texttt{Simplify}}

Mathematica nie wie automatycznie, że masa, długość albo liczba kwantowa są
dodatnie. Bez założeń może nie uprościć wyrażeń albo uprościć je w sposób trudny
do interpretacji. Dlatego należy używać:
\begin{verbatim}
Simplify[expr, Assumptions -> {a > 0, m > 0, n \[Element] Integers}]
\end{verbatim}
Dla całek kwantowych założenia są często konieczne. Dotyczy to szczególnie
parametrów takich jak \(L>0\), \(n\in\mathbb{N}\), \(m>0\), \(\hbar>0\).

\section{Nieodróżnianie przypadku \texorpdfstring{\(E<V_0\)}{E<V0} od \texorpdfstring{\(E>V_0\)}{E>V0}}

W barierze potencjału charakter rozwiązania zależy od znaku \(E-V_0\). Jeśli
\(E>V_0\), rozwiązanie w obszarze bariery jest oscylacyjne. Jeśli \(E<V_0\), jest
wykładniczo zanikające albo narastające.

Mylenie tych przypadków prowadzi do błędnych wzorów na transmisję i błędnych
wykresów. Przed rozpoczęciem rachunku trzeba zawsze napisać, w którym reżimie
pracujemy.

\section{Błędne użycie \texttt{=} zamiast \texttt{:=}}

W Mathematice znak \texttt{=} oblicza prawą stronę od razu, a \texttt{:=} oblicza ją
dopiero przy wywołaniu. Dla funkcji zależnych od parametrów i zmiennych zwykle
bezpieczniejsze jest \texttt{:=}:
\begin{verbatim}
f[x_] := x^2 + a
\end{verbatim}
Błąd wyboru definicji może prowadzić do sytuacji, w której zmiana parametru nie
zmienia wyniku, bo wyrażenie zostało obliczone wcześniej.

\section{Stare definicje w notebooku}

Notebook Mathematiki pamięta wcześniejsze definicje. Jeśli student zmienił nazwę
zmiennej albo usunął komórkę, stara definicja może dalej istnieć w jądrze.
Dlatego przed oddaniem projektu warto uruchomić notebook od początku po komendzie:
\begin{verbatim}
ClearAll["Global`*"]
\end{verbatim}
Jeśli notebook działa tylko po przypadkowej historii wcześniejszych uruchomień,
nie jest poprawnym projektem.

\section{Mylenie \texttt{E} i \texttt{I} w Mathematice}

W Mathematice \texttt{E} oznacza liczbę Eulera, a \texttt{I} oznacza jednostkę
urojoną. Nie należy używać \texttt{E} jako nazwy energii ani \texttt{I} jako indeksu.
Lepiej pisać na przykład:
\begin{verbatim}
energy[n_] := n^2 Pi^2 hbar^2/(2 m L^2)
\end{verbatim}
zamiast ryzykować konflikt z wbudowanymi symbolami.

\section{Brak \texttt{Evaluate} przy wykresach rozwiązań numerycznych}

Po użyciu \texttt{NDSolve} wynik często ma postać reguły podstawienia. Przy rysowaniu
rozwiązania przydatne jest \texttt{Evaluate}:
\begin{verbatim}
Plot[Evaluate[y[t] /. sol], {t, 0, tmax}]
\end{verbatim}
Bez tego wykres może działać wolniej, generować ostrzeżenia albo nie pokazywać
tego, czego student oczekuje.

\section{Zbyt brzydkie i nieczytelne wykresy}

Wykres jest częścią argumentu fizycznego. Powinien mieć podpisane osie, jednostki,
legendę i rozsądny zakres. Wykres bez jednostek jest prawie zawsze niekompletny.

Dobre minimum to:
\begin{itemize}
    \item opis osi z jednostkami,
    \item legenda, jeśli jest kilka krzywych,
    \item czytelny zakres,
    \item tytuł albo podpis mówiący, co przedstawia wykres,
    \item brak przypadkowych skal i niedopasowanych kolorów.
\end{itemize}

\section{Wynik numeryczny bez interpretacji}

Sama liczba albo wykres nie jest jeszcze odpowiedzią fizyczną. Trzeba napisać, co
oznacza wynik. Czy energia rośnie przy zmniejszaniu szerokości studni? Czy
transmisja maleje przy poszerzaniu bariery? Czy norma jest zachowana? Czy wynik
ma właściwy rząd wielkości?

Raport bez interpretacji wygląda jak zrzut z programu, a nie jak praca z fizyki.

\section{Kod, którego autor sam nie rozumie}

Najgorszy kod to kod skopiowany bez zrozumienia. W raporcie student powinien umieć
wyjaśnić każdą ważną linię: co oznaczają parametry, co liczy funkcja, jakie są
jednostki i dlaczego taki zakres wykresu został wybrany.

Jeśli kod jest długi, warto podzielić go na bloki: definicja stałych, definicja
funkcji, obliczenia, wykresy, kontrola wyniku. Kod powinien być narzędziem
myślenia, nie zasłoną dymną.

\section{Zbyt duże zaufanie do AI}

Narzędzia AI mogą pomóc w pisaniu kodu, szukaniu błędów i redagowaniu raportu,
ale nie gwarantują poprawności fizycznej. Student musi sprawdzić jednostki,
warunki brzegowe, normalizację i sens wyniku. Kod wygenerowany przez AI, którego
autor nie rozumie, jest tak samo problematyczny jak kod skopiowany z internetu.

\begin{summarybox}
Większość błędów w praktycznej mechanice kwantowej można wykryć przez pięć pytań:
czy jednostki są poprawne, czy funkcja jest znormalizowana, czy warunki brzegowe
są jasne, czy wykres jest opisany i czy wynik ma sens fizyczny.
\end{summarybox}
